% AUTO-GENERATED. DO NOT EDIT.
% Generator: hpc/paper_artifacts/build_tab_task_spec.py
% Benchmark: EvEMTBench v1.3.0
% Inputs: (none — static specification)
\begin{table*}[t]
\centering
%\small

\caption{Target and valid-sample definitions for the five primary EvEMTBench functions. L, Li, and G denote local, line, and global observability, respectively.}
\label{tab:task_spec}

\begin{tabular}{@{}
  >{\raggedright\arraybackslash}p{2.7cm}
  >{\raggedright\arraybackslash}p{3.0cm}
  >{\raggedright\arraybackslash}p{8.3cm}
  >{\centering\arraybackslash}p{1.5cm}
@{}}
\toprule
\textbf{Function} &
\textbf{Target} &
\textbf{Valid evaluation windows} &
\textbf{Views} \\
\midrule

Fault detection (\acs{fd})
& Fault active / inactive
& All windows from fault episodes; fault-active windows are positive, while windows before inception or after the annotated fault interval are negative
& L, Li, G \\

Fault classification (\acs{fc})
& Fault class
& Fault-active windows with labels in the frozen nine-class fault vocabulary
& L, Li, G \\

Fault localization (\acs{fl})
& Position along line (\% of line length)
& Fault-active windows from line-fault episodes; busbar faults are excluded because no unique line position is defined
& L, Li, G \\

Event detection (\acs{ed})
& Event active / inactive
& All valid event-analysis windows; fault- or switching-active windows are positive, while inactive windows are negative
& G \\

Event classification (\acs{ec})
& Event class
& All valid event-analysis windows over the frozen 23-class vocabulary; inactive windows are labeled \texttt{no\_event}
& G \\

\bottomrule
\end{tabular}
\end{table*}
